% ============================================================
% DS MOD08 -- SUJET A : loi binomiale, loi de Poisson, loi normale, dénombrement
% BTS ATI / CRSA 1re année -- durée 1 h 30
% Pôle Formation UIMM - CVDL - S. JAUBERT
% Corrigé et barème : MOD08_DS_Binomiale_Poisson_Normale_ATI_CRSA_Sujet_A_Corrige_v1_0 (version 4 parties : le sujet B, en 3 parties, n'est plus utilisé)
% ============================================================
\documentclass[a4paper,11pt]{article}
\usepackage[utf8]{inputenc}
\usepackage[T1]{fontenc}
% babel-french absent sur certains postes : chargement conditionnel
\IfFileExists{french.ldf}{\usepackage[french]{babel}}{}
\usepackage{helvet}
\usepackage{courier}
\renewcommand{\familydefault}{\sfdefault}
\usepackage[top=3cm, bottom=2.5cm, left=2cm, right=2cm, headheight=2cm]{geometry}
\usepackage{amsmath,amssymb,amsfonts}
\usepackage{graphicx}
\usepackage[dvipsnames,table]{xcolor}
\usepackage{tikz}
\usepackage[most]{tcolorbox}
\usepackage{fancyhdr}
\usepackage{enumitem}
\usepackage{array,needspace}
\usepackage{titlesec}
\usepackage{hyperref}

\definecolor{uimmblue}{HTML}{005C85}
\definecolor{uimmred}{HTML}{D31326}
\definecolor{uimmlightblue}{HTML}{E8F1F8}
\definecolor{uimmgray}{HTML}{333333}
\definecolor{uimmlightgray}{HTML}{F5F5F5}
\hypersetup{colorlinks=true, linkcolor=uimmblue, urlcolor=uimmblue}

\titleformat{\section}{\Large\bfseries\color{uimmblue}}{}{0pt}{}[\vspace{2pt}{\color{uimmblue}\titlerule[1.2pt]}]
\titleformat{\subsection}{\large\bfseries\color{uimmblue}}{}{0pt}{}
\setlist[enumerate,1]{label={\color{uimmblue}\textbf{\arabic*.}}, leftmargin=*}
\setlist[enumerate,2]{label=\alph*), leftmargin=*}
\setlist[itemize,1]{label={\color{uimmblue}\textbullet}}

\newtcolorbox{exo}[1]{colback=white, colframe=uimmblue, fonttitle=\bfseries\color{white},
  coltitle=white, title={#1}, boxrule=1pt, arc=2pt,
  left=8pt, right=8pt, top=5pt, bottom=6pt, before skip=10pt, after skip=6pt}
\newtcolorbox{consigne}{colback=uimmlightgray, colframe=uimmgray, boxrule=0.8pt, arc=3pt,
  left=8pt, right=8pt, top=6pt, bottom=6pt}

\newcommand{\pts}[1]{\hfill\mbox{\color{uimmred}\small\textbf{(#1)}}}
% Zone de réponse lignée : #1 = nombre de lignes
\newcommand{\zone}[1]{\par\smallskip\noindent\begin{tikzpicture}
  \foreach \i in {1,...,#1}{\draw[uimmgray!30] (0,-0.85*\i) -- (\linewidth-1pt,-0.85*\i);}
  \end{tikzpicture}\par\medskip}
\newcommand{\EnTete}[2]{%
\pagestyle{fancy}\fancyhf{}\renewcommand{\headrulewidth}{0pt}
\fancyhead[L]{\raisebox{-0.5\height}{\includegraphics[height=1.5cm]{logo_uimm_placeholder.jpg}}}
\fancyhead[C]{\begin{minipage}{7cm}\centering
  {\small\color{uimmblue}\textbf{Pôle Formation UIMM - CVDL}}\\
  {\footnotesize\color{uimmgray} BTS ATI / CRSA 1\textsuperscript{re} année\\ Mathématiques}\end{minipage}}
\fancyhead[R]{{\small\color{uimmred}\textbf{MOD08 -- #1}}}
\fancyfoot[C]{{\color{uimmblue}\rule{\textwidth}{1pt}}\\[2pt]
  {\footnotesize\color{uimmgray} Pôle Formation UIMM -- Centre-Val de Loire \hfill S. JAUBERT \hfill #2 -- \thepage}}}
\EnTete{DS Sujet A}{Sujet A -- DS 1 h 30}
\newtcolorbox{rappel}{colback=uimmlightblue, colframe=uimmblue, fonttitle=\bfseries\color{white},
  coltitle=white, title={Rappels de cours}, boxrule=1pt, arc=2pt,
  left=8pt, right=8pt, top=5pt, bottom=6pt, before skip=8pt, after skip=8pt}

\begin{document}

\begin{tcolorbox}[colback=uimmred, colframe=uimmred, sharp corners, boxrule=0pt,
  left=14pt, right=14pt, top=10pt, bottom=10pt]
{\color{white}\LARGE\bfseries Devoir surveillé -- MOD08 \hfill Sujet A}\\[4pt]
{\color{white}\large Lois binomiale, de Poisson et normale~; dénombrement \hfill Durée~: 1 h 30}
\end{tcolorbox}

\vspace{4pt}
\noindent{\small\renewcommand{\arraystretch}{1.6}%
\begin{tabular}{@{}>{\bfseries\color{uimmblue}}l@{\hspace{6pt}}p{8.6cm}@{\hspace{10pt}}>{\bfseries\color{uimmblue}}l@{\hspace{6pt}}p{3.2cm}@{}}
Nom, prénom~: & \dotfill & Groupe~: & \dotfill \\
 & & Note~: & \dotfill\ / 20 \\
\end{tabular}}

\begin{consigne}
\textbf{Consignes.} Calculatrice autorisée en mode examen, aucun document. Répondre sur le sujet, dans les zones prévues.
\begin{itemize}[leftmargin=12pt, itemsep=1pt, topsep=2pt]
  \item Toute loi utilisée doit être \textbf{nommée avec ses paramètres}.
  \item Pour chaque probabilité demandée~: écrire l'évènement à l'aide de la variable aléatoire, \textbf{puis} le résultat. Un résultat seul ne rapporte pas tous les points.
  \item Probabilités arrondies à $10^{-4}$. Le barème est donné à titre indicatif~: partie 1 sur 6 points, partie 2 sur 6 points, partie 3 sur 4 points, partie 4 sur 4 points.
\end{itemize}
\end{consigne}

\begin{rappel}
\textbf{Loi binomiale $\mathcal{B}(n\,;p)$.} On répète $n$ fois, de façon identique et indépendante, une épreuve à deux issues (succès de probabilité $p$). $X$ compte le nombre de succès~; elle prend les valeurs $0, 1, \dots, n$.
\[
P(X = k) = \binom{n}{k}\,p^{k}\,(1-p)^{n-k},\quad \binom{n}{k} = \frac{n!}{k!\,(n-k)!},\quad E(X) = np,\quad \sigma(X) = \sqrt{np(1-p)}.
\]
\textbf{Loi de Poisson $\mathcal{P}(\lambda)$.} Nombre d'évènements rares (pannes, défauts, alarmes) sur une durée ou une zone donnée. Valeurs $0, 1, 2, \dots$
\[
P(X = k) = e^{-\lambda}\,\frac{\lambda^{k}}{k!},\qquad E(X) = \lambda,\qquad \sigma(X) = \sqrt{\lambda}.
\]
\textbf{Approximation de $\mathcal{B}(n\,;p)$ par $\mathcal{P}(\lambda)$, avec $\lambda = np$.} Elle est légitime lorsque les \textbf{trois} conditions sont réunies~: \ $n \geq 30$, \ $p \leq 0{,}1$, \ $np \leq 10$.

\smallskip
\textbf{Loi normale $\mathcal{N}(\mu\,;\sigma)$.} $\mu$ est la moyenne, $\sigma$ l'\textbf{écart-type}. Règle des 68 -- 95 -- 99,7~:\\
$P(\mu-\sigma\leq X\leq\mu+\sigma)\approx 0{,}68$, \ $P(\mu-2\sigma\leq X\leq\mu+2\sigma)\approx 0{,}95$, \ $P(\mu-3\sigma\leq X\leq\mu+3\sigma)\approx 0{,}997$.

\smallskip
\textbf{Dénombrement.} Ordonner $n$ objets distincts (permutations)~: $n! = n\times(n-1)\times\dots\times 2\times 1$ façons. Choisir $k$ objets parmi $n$, sans tenir compte de l'ordre (combinaisons)~: $\binom{n}{k}$ façons. Si toutes les issues ont la même probabilité~:
\[
P(A) = \frac{\text{nombre de cas favorables à } A}{\text{nombre de cas possibles}}.
\]

\smallskip
\textbf{Calculatrice.} «~Fdp~» (ou «~pd~») donne $P(X = k)$~; «~FRép~» (ou «~cd~») donne $P(X \leq k)$.\\
{\small TI~: \texttt{binomFdp(n,p,k)}, \texttt{binomFRép(n,p,k)}, \texttt{poissonFdp($\lambda$,k)}, \texttt{poissonFRép($\lambda$,k)}, \texttt{normalFRép(a,b,$\mu$,$\sigma$)}.\\
Casio~: \texttt{Bpd}, \texttt{Bcd}, \texttt{Ppd}, \texttt{Pcd}, \texttt{Ncd}. \quad Borne infinie pour la loi normale~: \texttt{-1E99} ou \texttt{1E99}.\\
Dénombrement (menu \texttt{PROB})~: \texttt{!} (factorielle), \texttt{nCr} (combinaisons).}
\end{rappel}

% ============================================================
\newpage
\section{Partie 1 -- Loi binomiale \hfill {\normalsize\color{uimmred}6 points}}
\begin{exo}{Réception de détecteurs inductifs}
Un intégrateur de cellules robotisées achète des détecteurs de proximité inductifs M12 par cartons de plusieurs centaines. Le fournisseur annonce que $5\,\%$ des détecteurs présentent un défaut (distance de commutation hors tolérance). À la réception, le technicien prélève $25$ détecteurs et les teste sur un banc. L'arrivage est assez important pour assimiler ce prélèvement à un tirage avec remise.\\
On note $X$ le nombre de détecteurs \textbf{défectueux} dans le prélèvement.
\end{exo}
\begin{enumerate}
  \needspace{7.4cm}
  \item Justifier que $X$ suit une loi binomiale. Préciser ses paramètres. \pts{1 pt}
  \zone{6}
  \needspace{6.5cm}
  \item Calculer l'espérance $E(X)$ et l'écart-type $\sigma(X)$. Interpréter $E(X)$ dans le contexte. \pts{1 pt}
  \zone{5}
  \needspace{6.5cm}
  \item Calculer la probabilité que le prélèvement contienne exactement $2$ détecteurs défectueux. Détailler le calcul à l'aide de la formule, \textbf{sans} utiliser les fonctions de loi de la calculatrice. \pts{1 pt}
  \zone{5}
  \needspace{4cm}
  \item Pour chacun des évènements suivants, traduire l'évènement à l'aide de $X$, puis calculer sa probabilité.
  \begin{enumerate}
    \needspace{4cm}
    \item Aucun détecteur du prélèvement n'est défectueux. \pts{0,5 pt}
    \zone{3}
    \needspace{4.8cm}
    \item Le prélèvement contient au moins un détecteur défectueux. \pts{0,5 pt}
    \zone{4}
    \needspace{5.6cm}
    \item Au moins $23$ détecteurs du prélèvement fonctionnent correctement. \pts{0,75 pt}
    \zone{5}
    \needspace{4.8cm}
    \item Le prélèvement contient au plus $3$ détecteurs défectueux. \pts{0,75 pt}
    \zone{4}
  \end{enumerate}
  \needspace{5.6cm}
  \item Règle de réception~: l'arrivage est \textbf{refusé} si le prélèvement contient plus de $3$ détecteurs défectueux. Calculer la probabilité de refuser un arrivage dont la proportion de défectueux est bien de $5\,\%$. \pts{0,5 pt}
  \zone{4}
\end{enumerate}

% ============================================================
\newpage
\section{Partie 2 -- Loi de Poisson \hfill {\normalsize\color{uimmred}6 points}}
\subsection*{A. Pertes de communication sur le réseau de terrain}
\begin{exo}{Contexte}
Sur une ligne d'assemblage, l'automate dialogue avec des îlots d'entrées/sorties déportés par un réseau Profinet. Le diagnostic de l'automate enregistre des \textbf{pertes de communication} brèves, qui surviennent de façon aléatoire et indépendante. L'historique montre qu'il se produit \textbf{en moyenne $2{,}5$ pertes de communication par semaine} de production. On admet que le nombre de pertes au cours d'une semaine suit une loi de Poisson.
\end{exo}
\begin{enumerate}
  \needspace{5.6cm}
  \item Définir une variable aléatoire $Y$ adaptée à la situation. Donner sa loi et justifier la valeur de son paramètre. \pts{0,75 pt}
  \zone{3}
  \needspace{4cm}
  \item Calculer l'écart-type de $Y$. \pts{0,25 pt}
  \zone{2}
  \needspace{6.5cm}
  \item Calculer la probabilité qu'il ne se produise aucune perte de communication dans la semaine, puis la probabilité qu'il s'en produise exactement $3$. Détailler les calculs à l'aide de la formule. \pts{0,75 pt}
  \zone{4}
  \needspace{4.8cm}
  \item Calculer la probabilité qu'il se produise plus de $4$ pertes dans la semaine. \pts{0,75 pt}
  \zone{4}
  \needspace{5.4cm}
  \item Le technicien réseau veut fixer un seuil $k$~: si une semaine compte plus de $k$ pertes, il lance un diagnostic complet du réseau. Déterminer la plus petite valeur de $k$ telle que $P(Y \leq k) \geq 0{,}95$. Justifier par deux probabilités. \pts{1 pt}
  \zone{4}
\end{enumerate}

\newpage
\subsection*{B. Approximation d'une loi binomiale par une loi de Poisson}
\begin{exo}{Contexte}
Un fabricant livre des cartes d'entrées TOR par lots de $400$ cartes. La probabilité qu'une carte soit défectueuse est $0{,}0075$. Le lot est assimilé à un tirage avec remise. On note $N$ le nombre de cartes défectueuses dans un lot.
\end{exo}
\begin{enumerate}[resume]
  \needspace{4cm}
  \item Donner la loi exacte de $N$ et ses paramètres. \pts{0,5 pt}
  \zone{3}
  \needspace{5.6cm}
  \item Justifier, en vérifiant chacune des conditions, que l'on peut approcher la loi de $N$ par une loi de Poisson. Préciser son paramètre. \pts{1 pt}
  \zone{6}
  \needspace{4.8cm}
  \item À l'aide de cette loi de Poisson, calculer la probabilité qu'un lot contienne au plus $2$ cartes défectueuses. \pts{0,5 pt}
  \zone{5}
  \needspace{4.8cm}
  \item Calculer la même probabilité avec la loi exacte. Comparer les deux résultats et conclure sur la qualité de l'approximation. \pts{0,5 pt}
  \zone{5}
\end{enumerate}

% ============================================================
\newpage
\section{Partie 3 -- Loi normale \hfill {\normalsize\color{uimmred}4 points}}
\begin{exo}{Temps de cycle d'un robot de palettisation}
En fin de ligne, un robot de palettisation prélève des cartons sur un convoyeur et les dépose sur une palette. On note $T$ le temps de cycle du robot, en secondes, pour un carton pris au hasard. On admet que $T$ suit la loi normale $\mathcal{N}(12{,}0\,;0{,}4)$.
\end{exo}
\begin{enumerate}
  \needspace{4.8cm}
  \item Calculer la probabilité qu'un temps de cycle soit compris entre $11{,}5$\,s et $12{,}5$\,s. \pts{1 pt}
  \zone{4}
  \needspace{4.8cm}
  \item La cadence de la ligne impose un temps de cycle d'au plus $12{,}8$\,s. Calculer la probabilité qu'un cycle dépasse $12{,}8$\,s. \pts{1 pt}
  \zone{4}
  \needspace{5.6cm}
  \item À l'aide de la règle des 68 -- 95 -- 99,7, donner l'intervalle qui contient environ $95\,\%$ des temps de cycle. Ce résultat est-il cohérent avec celui de la question 2~? Justifier. \pts{1 pt}
  \zone{5}
  \needspace{4.8cm}
  \item Au cours d'un poste, le robot réalise $1\,000$ cycles. Estimer le nombre moyen de cycles qui dépassent $12{,}8$\,s. \pts{1 pt}
  \zone{4}
\end{enumerate}

% ============================================================
\newpage
\section{Partie 4 -- Dénombrement \hfill {\normalsize\color{uimmred}4 points}}
\begin{exo}{Vissage robotisé d'un carter}
Dans une cellule d'assemblage, un robot équipé d'une visseuse électrique serre les $6$ vis de fixation d'un carter de motoréducteur. Les vis sont repérées $V_1$, $V_2$, \dots, $V_6$. L'ordre dans lequel le robot serre les vis est défini dans son programme.
\end{exo}
\begin{enumerate}
  \needspace{4.8cm}
  \item Combien d'ordres de vissage différents le programmeur peut-il définir~? \pts{1 pt}
  \zone{4}
  \needspace{4.8cm}
  \item Pour garantir la mise en position du carter, la vis $V_1$ doit être serrée en premier. Combien d'ordres de vissage restent possibles~? \pts{0,5 pt}
  \zone{4}
  \needspace{5.6cm}
  \item En fin de poste, le contrôleur qualité vérifie à la clé dynamométrique le couple de serrage de $3$ vis du carter. L'ordre du contrôle n'a pas d'importance. Combien de groupes de $3$ vis différents peut-il contrôler~? \pts{0,75 pt}
  \zone{4}
  \needspace{4cm}
  \item Sur un carter, $2$ des $6$ vis sont mal serrées (couple insuffisant).
  \begin{enumerate}
    \needspace{4cm}
    \item Combien de groupes de $3$ vis ne contiennent aucune vis mal serrée~? \pts{0,5 pt}
    \zone{3}
    \needspace{4.8cm}
    \item Combien de groupes de $3$ vis contiennent exactement une vis mal serrée~? \pts{0,5 pt}
    \zone{4}
  \end{enumerate}
  \needspace{5.6cm}
  \item Le contrôleur choisit les $3$ vis au hasard~: tous les groupes de $3$ vis ont la même probabilité d'être choisis. Calculer la probabilité que le contrôle détecte au moins une vis mal serrée. \pts{0,75 pt}
  \zone{5}
\end{enumerate}

\end{document}
